\subsection{Relación objetivos, fases y procedimiento}
\label{sec:objetivos-metodologia}

La \Cref{tab:oe-metodologia} pone cada OE en una fase, con un verbo y un producto.

\begin{table}[H]
  \centering
  \caption{Objetivos específicos vinculados al procedimiento metodológico}
  \label{tab:oe-metodologia}
  \small
  \begin{tabular}{@{}p{0.06\textwidth}p{0.14\textwidth}p{0.22\textwidth}p{0.24\textwidth}p{0.24\textwidth}@{}}
    \toprule
    \textbf{OE} & \textbf{Fase} & \textbf{Acción (verbo)} & \textbf{Entrada} & \textbf{Producto} \\
    \midrule
    --- &
    Fase 1 &
    \textbf{Diseñar} entradas, scores y umbrales &
    Requisitos del §1.4; Anexo C &
    Modelo S/E/C documentado \\
    \textbf{OE1} &
    Fase 2 &
    \textbf{Adquirir} y \textbf{normalizar} contrato y SD &
    OpenAPI + extracto BIAN (§4.3) &
    Representaciones comparables \\
    \textbf{OE2} &
    Fase 3 &
    \textbf{Emparejar} estructuras y calcular E &
    Representaciones OE1 &
    Matriz de correspondencias + E \\
    \textbf{OE3} &
    Fase 4 &
    \textbf{Calcular} similitud semántica S &
    Representaciones OE1 &
    Matriz semántica + S \\
    \textbf{OE4} &
    Fase 5 &
    \textbf{Integrar} S, E, C y clasificar &
    E, S, representaciones OE1 &
    C, \AlignmentScore{}, tipologías \\
    \bottomrule
  \end{tabular}
\end{table}

Las técnicas por OE están en la \Cref{sec:tecnicas-herramientas}. Las filas, en el Anexo~C.
